这本讲义换一种方式组织数学物理方法：不从数学分支出发，而从物理问题出发。市面上的教材多数按「复变函数—数学物理方程—特殊函数」排列。这个顺序对写书的人是自然的，对学物理的人是反的：学生先学完一整本复分析，还不知道这些定理将来解的是哪道题。另一种常见的毛病是讲得太细、太严格——洛朗级数的收敛性证明占据半节课，而学生连为什么需要这种展开都没弄明白。这种严格在物理里收益很低，代价却很大：它挤掉了真正需要讲的东西，也就是一个工具什么时候能用、什么时候会坏。所以正文共十章（另设第 0 章导论）。每一章围绕一个具体的物理问题展开：一根弦怎么振动、一面鼓为什么只发出某些频率、点电荷在导体腔里的电势怎么算、机翼的升力从哪来、脉冲进入系统后响应为什么不可能提前出现。数学工具在这些问题卡住的地方出场，出场时必然已经有一个问题在等它。与上一版相比，这一版做了几件具体的事。第一，推导写细了：凡是跳过一步代数，读者就会卡在那里，所以本书尽量把每一步的理由写出来，包括用到的近似丢掉了哪一阶。第二，题多了：每章三十余道计算题、近十道概念判断题，全部配详解，且尽量取自真实的物理情形，不凑偏难怪题。第三，图多了：本征模、节线、回路、镜像、映射、波包展宽这些「光看公式看不出来」的东西，都画出来了。第四，章拆小了：圆形的振动（贝塞尔函数）与球形的世界（勒让德函数）各自独立成章，复平面上的留数与保角映射也各自独立成章，每一章都能在两三次课内读完。还有一条约定值得单独说明：本书**把偏导数写全**，并把拉普拉斯算子统一写作 $\nabla^2$。$\partial^2 u/\partial x^2$ 比 $u_{xx}$ 长，但读者不需要在脑子里做符号翻译：你可以在任何时候把式子念出来，念出来就是它本来的意思。同理，$\nabla^2$ 在全书中只有一种写法，不用正三角 $\Delta$，也不与其他记号混用——符号应当节省读者的力气，而不是节省作者的排版。本书不追求数学体系的完整，但必须保住「什么时候会坏」的判断能力。这是它和「只讲套路」的教材最根本的区别。每章末尾有一份判断点清单，回答三个问题：这个操作什么时候合法、假设破了会怎样、怎么在具体问题里察觉。附录 B 把全书判断点汇总成一张速查表。本书由胡甜与 DeepSeek 协作完成。选题与讲法参考了 Arfken、Boas、Riley、Butkov、Courant–希尔伯特、Jackson、顾樵等教材的**思路**，正文的组织、推导、例题与习题解答均为独立撰写。参考书目见书末。
